% Source: https://artofproblemsolving.com/community/c3h1331017_helpe_me_plz
% Stable identity: aops-topic-1331017-q01
Let $u=\sqrt{x+1}\ge1$ and $v=\sqrt{y+2}\ge\sqrt2$. The condition becomes $u^2-1-3u=3v-v^2+2$, i.e.
\[
u^2+v^2=3(u+v)+3,\qquad\text{equivalently}\qquad\left(u-\dfrac{3}{2}\right)^2+\left(v-\dfrac{3}{2}\right)^2=\dfrac{15}{2}.
\]
Put $s=u+v$. Then $x+y=u^2+v^2-3=3s$ and $uv=\dfrac{s^2-(u^2+v^2)}{2}=\dfrac{s^2-3s-3}{2}$, so
\[
S=3s-uv=\dfrac{-s^2+9s+3}{2}.
\]
The point $(u,v)$ runs over the arc of the circle with $u\ge1$, $v\ge\sqrt2$. On this arc, $s=u+v$ is largest at $u=v=\dfrac{3+\sqrt{15}}{2}$, where $s=3+\sqrt{15}$, and smallest at an end of the arc. At the end $u=1$ we get $v=\dfrac{3+\sqrt{29}}{2}$ and $s=\dfrac{5+\sqrt{29}}{2}\approx5.19$; at the end $v=\sqrt2$, $u\approx4.24$ and $s\approx5.65$. So $s$ ranges over $\left[\dfrac{5+\sqrt{29}}{2},\,3+\sqrt{15}\right]$, which lies to the right of $s=\dfrac{9}{2}$, where the quadratic $S(s)$ peaks; hence $S$ is decreasing in $s$ on this range.
\begin{itemize}
\item Maximum, at $s=\dfrac{5+\sqrt{29}}{2}$ (that is, $x=0$, $y=\dfrac{15+3\sqrt{29}}{2}$): $S=\dfrac{-s^2+9s+3}{2}=\boxed{6+\sqrt{29}}$.
\item Minimum, at $s=3+\sqrt{15}$ (that is, $x=\dfrac{11+3\sqrt{15}}{2}$, $y=\dfrac{9+3\sqrt{15}}{2}$): $S=\boxed{\dfrac{6+3\sqrt{15}}{2}}$.
\end{itemize}
